
\subsubsection{5.1 Main Finding}

The CV-based detection mechanism achieves AUC = 0.0, decisively failing the $\geq$ 0.75 gate.

\subsubsection{5.2 CV Values}

\begin{table}[htbp]
\centering
\resizebox{\columnwidth}{!}{%
\begin{tabular}{lll}
\toprule
Feature Type & CV Value & Expected Direction \\
\midrule
Background (spurious) & 0.0393 & Low (< 0.15) \\
Bird Type (core) & 0.0360 & High (> 0.20) \\
Difference & 0.0033 & --- \\
\bottomrule
\end{tabular}
}
\end{table}

Both feature types exhibit nearly identical CV values (~0.04). The hypothesized separation does not manifest.

\subsubsection{5.3 Direction Reversal}

The hypothesis predicted CV(spurious) < CV(core). The observed values show CV(spurious) = 0.0393 > CV(core) = 0.0360, a marginal reversal of the expected direction. With only two feature types, this reversal yields AUC = 0.0.

\subsubsection{5.4 Gate Evaluation}

\begin{table}[htbp]
\centering
\resizebox{\columnwidth}{!}{%
\begin{tabular}{llll}
\toprule
Metric & Value & Threshold & Status \\
\midrule
AUC & 0.0 & $\geq$ 0.75 & FAIL \\
Best F1 & 0.667 & --- & --- \\
\bottomrule
\end{tabular}
}
\end{table}

\subsubsection{5.5 Probe Trajectory Analysis}

Both background and bird type probes show flat trajectories across the C sweep. Accuracy is near 100\% at all regularization strengths for both feature types. There is no observable emergence pattern---features are already fully learned in the CLIP representation.

